\chapter{Anexo E: operacion, API y despliegue}

\section{API}

\path{doc/docs/operations/api-deployment.md} documenta la API de prediccion y explicabilidad local. El servicio carga modelos, aplica el mismo procesamiento de features y devuelve predicciones con metadatos. La explicabilidad local se apoya en los artefactos del modelo seleccionado y en la estructura de dashboard.

\section{Almacenamiento}

La configuracion permite almacenamiento local y GCP. En local, los modelos entrenados se buscan en \path{src/volatility_models/trained_models}, los metadatos en \path{src/volatility_models/trained_metadata/retrained_metadata} y los bundles en \path{src/dashboard/saved_models}. En GCP se parametrizan buckets y prefijos.

\section{Docker y Terraform}

Los dockerfiles empaquetan servicios. Terraform define infraestructura cuando se use despliegue cloud. Aunque la evaluacion cuantitativa no depende de estas piezas, su presencia indica que el proyecto se ha pensado como sistema desplegable.

\section{Activos estaticos de documentacion}

La carpeta \path{doc/docs/javascripts/} contiene configuracion de MathJax, inicializacion de Mermaid y ajustes de etiquetas de navegacion. \path{doc/docs/stylesheets/extra.css} contiene estilos adicionales. Estos ficheros no alteran la metodologia financiera, pero forman parte de la documentacion navegable y explican por que los diagramas y formulas se muestran correctamente en MkDocs.
